\section{Review of Macdonald polynomials}\label{sec:macdonald_polynomials}

Here we collect the necessary notation and results
around Macdonald symmetric polynomials.
We follow \cite[Chapter VI]{Macdonald1995}.

\subsection{Definition}\label{sub:Macdonald_polynomials_definition}

Let $N\ge1$.
Macdonald symmetric polynomials $P_\lambda$ in $N$ variables $x_1,\dots,x_N$
are indexed by partitions
$\lambda=(\lambda_1\ge \dots\ge \lambda_N \ge0)$, $\lambda_i\in \mathbb{Z}$,
with $N$ parts. Denote the set of these partitions by~$\mathbb{Y}(N)$.
The $P_\lambda$'s
depend on two parameters $q,t\in[0,1)$. For each fixed
$(q,t)$, they form a~basis in the space of symmetric
polynomials in~$N$ variables when $\lambda$ runs over~$\mathbb{Y}(N)$.
The shortest definition of the
$P_\lambda$'s is through the first Macdonald $q$-difference operator
acting in the~$x_i$'s:
\begin{gather*}
	D_1\coloneqq \sum_{i=1}^{N}
	\biggl( \prod_{1\le j\le N\colon j\ne i}\frac{x_j-tx_i}{x_j-x_i} \biggr)
	T_{q;i},\\ T_{q;i}f(x_1,\dots,x_N )\coloneqq f(x_1,\dots,x_{i-1},qx_i,x_{i+1},\dots,x_N ).
\end{gather*}
The operator $D_1$ preserves the space of symmetric
polynomials in $x_1,\dots,x_N $, and its eigenfunctions
are the Macdonald polynomials
\begin{gather}
	D_1P_\lambda(x_1,\dots,x_N\mid q,t )=
	\bigl( q^{\lambda_1}t^{N-1}+q^{\lambda_2}t^{N-2}+\dots+q^{\lambda_N}
	\bigr)
	P_\lambda(x_1,\dots,x_N\mid q,t ),\label{eq:Macdonald_eigenrelation}
\end{gather}
with $\lambda\in \mathbb{Y}(N)$.
For generic $(q,t)$, the eigenvalues in~\eqref{eq:Macdonald_eigenrelation}
for different $\lambda$ are distinct, so
the $P_\lambda$'s are determined uniquely up to normalization.
The normalization is specified by
\begin{equation*}
	P_\lambda(x_1,\dots,x_N\mid q,t )=
	x_1^{\lambda_1}x_2^{\lambda_2}\cdots x_N^{\lambda_N} +
	\text{lexicographically lower terms},
\end{equation*}
where the lower terms depend on $q$, $t$.

In the case $q=t$, the polynomials $P_\lambda$ reduce to the
well-known Schur symmetric polynomials~$s_\lambda$, which admit
the following explicit determinantal formula (which does not
depend on the choice of $q$):
\begin{equation}
	\label{eq:Schur}
	P_\lambda(x_1,\dots,x_N\mid q,q )=s_\lambda(x_1,\dots,x_N )=
	\frac{\det\bigl[ x_i^{\lambda_j+N-j} \bigr]_{i,j=1}^{N}}{\det
	\bigl[ x_i^{N-j} \bigr]_{i,j=1}^N}.
\end{equation}
For generic $(q,t)$, there are no known formulas for Macdonald polynomials which are as compact as~\eqref{eq:Schur}.

\subsection{Principal specialization}\label{sub:principal_specialization}

\begin{figure}[htpb]
	\centering
	\includegraphics[width=.28\textwidth]{fig_Young_d.pdf}
	\caption{Young diagram $\lambda=(5,5,5,2,1,1)$. Highlighted
		are the arm, leg, coarm, and coleg of the box $\square=(3,2)$.
		We have $|\lambda|=19$ and
		$a(\square)=3$, $l(\square)=1$, $a'(\square)=1$, $l'(\square)=2$.}
	\label{fig:Young_d}
\end{figure}

When the variables $x_i$ are specialized to a finite geometric
progression with the ratio $t$ (the second Macdonald parameter),
the polynomials $P_\lambda$ admit an explicit product formula in terms
of the Young diagram corresponding to the partition $\lambda$.
Recall that the Young diagram $\lambda$ is a~collection of $1\times 1$ boxes in the plane
with $\lambda_i$ boxes in row $i$, see Figure~\ref{fig:Young_d}
for an illustration.
The principal specialization takes the form
\cite[formulas~(VI.6.11) and (VI.6.11$'$)]{Macdonald1995}:
\begin{equation}	\label{eq:principal_spec}
	P_\lambda\big(1,t,\dots,t^{N-1} \mid q,t \big)=
	t^{n(\lambda)}\prod_{\square\in \lambda}
	\frac{1-q^{a'(\square)}t^{N-l'(\square)}}
	{1-q^{a(\square)}t^{l(\square)+1}},
\end{equation}
where the product is over all boxes of the Young diagram $\lambda$,
\begin{equation}	\label{eq:n_lambda_notation}
	n(\lambda):=\sum_i(i-1)\lambda_i,
\end{equation}
and the arms, legs, coarms, colegs of the box $\square=(i,j)$
are defined, respectively, as
\begin{equation}	\label{eq:arms_legs}
	a(\square)=\lambda_i-j,\qquad l(\square)=\lambda_j'-i,\qquad
	a'(\square)=j-1,\qquad l'(\square)=i-1.
\end{equation}
Here $\lambda_j'$ are the column lengths in the Young
diagram, see Figure~\ref{fig:Young_d}.

\subsection{Branching}\label{sub:branching}

Let us first recall the Pieri coefficients for Macdonald polynomials
\cite[formula~(VI.6.24.ii)]{Macdonald1995},
\cite[formula~(2.11)]{BorodinCorwin2011Macdonald}.
They depend on a pair of partitions $\mu$, $\lambda$
which \emph{interlace}, namely,
\begin{equation*}%\label{eq:interlace}
	\lambda_1\ge \mu_1\ge \lambda_2\ge \mu_2\ge \cdots
\end{equation*}
(notation $\mu\prec \lambda$), and are defined as
\begin{equation}
	\label{eq:psi_Pieri}
	\psi_{\lambda/\mu}
	=
	\psi_{\lambda/\mu}(q,t)
	\coloneqq
	\prod_{1\le i<j\le \ell(\mu)}
	\frac{\mathsf{f}\big(q^{\mu_i-\mu_j}t^{j-i}\big)
	\hspace{1pt}\mathsf{f}\big(q^{\lambda_i-\lambda_{j+1}}t^{j-i}\big)}
	{\mathsf{f}\big(q^{\lambda_i-\mu_j}t^{j-i}\big)
	\hspace{1pt}\mathsf{f}\big(q^{\mu_i-\lambda_{j+1}}t^{j-i}\big)},
	\qquad
	\mathsf{f}(u):=\frac{(tu;q)_{\infty}}{(qu;q)_{\infty}}.
\end{equation}
Here and below by $\ell(\mu)$ we denote the number of
parts in $\mu$ which are strictly positive,
that is, $\mu_{\ell(\mu)}>0$, $\mu_{\ell(\mu)+1}=0$ (when $\mu=(0,0,\dots )$, we set $\ell(\mu)=0$).
Let also $|\lambda|=\lambda_1+\dots+\lambda_{\ell(\lambda)}$
denote the number of boxes in the Young diagram~$\lambda$.

\begin{Proposition}[{\cite[formula~(VI.7.13$'$)]{Macdonald1995}}]
	\label{prop:branching}
	Let $\lambda\in \mathbb{Y}(N)$.
	We have
	\begin{equation}
		\label{eq:branching}
		P_\lambda(x_1,\dots,x_N \mid q,t)
		=
		\sum_{\mu\colon\mu\prec\lambda}
		\psi_{\lambda/\mu}(q,t)\hspace{1pt} P_\mu(x_1,\dots,x_{N-1} \mid q,t )\hspace{1pt}
		x_N^{|\lambda|-|\mu|},
	\end{equation}
	where the sum is over
	$\mu\in \mathbb{Y}(N-1)$ which interlace with $\lambda$.
\end{Proposition}

Note that for $q,t\in [0,1)$ the coefficients $\psi_{\lambda/\mu}$
\eqref{eq:psi_Pieri} are all nonnegative.
Together with Proposition~\ref{prop:branching} this implies:
\begin{Corollary}
	\label{cor:nonnegative_spec}
	Let $q,t\in[0,1)$. Specializing the variables
	$x_1,\dots,x_N $ into
	nonnegative real numbers
	makes the Macdonald polynomial
	$P_\lambda(x_1,\dots,x_N\mid q,t)$ nonnegative.
\end{Corollary}

More generally, for $N\ge K\ge1$
define the skew Macdonald polynomials
$P_{\lambda/\mu}$
as the coefficients
in the expansion:
\begin{equation*}
	P_\lambda(x_1,\dots,x_N \mid q,t)=
	\sum_{\mu\in \mathbb{Y}(K)}
	P_\mu(x_1,\dots,x_K \mid q,t )
	P_{\lambda/\mu}(x_{K+1},\dots,x_{N-1},x_N \mid q,t ).
\end{equation*}
Here we use the fact that
the polynomials
$P_\mu(x_1,\dots,x_K )$
form a basis in the space of symmetric functions in $K$ variables, and
expand
$P_\lambda(x_1,\dots,x_N )$
in this basis.
The Pieri coefficient is related to the skew Macdonald polynomial
in one variable as follows:
\begin{equation*}%\label{eq:P_skew_single_variable}
	P_{\lambda/\mu}(x_1\mid q,t)=
	\begin{cases}
		\psi_{\lambda/\mu}(q,t)\hspace{1pt} x_1^{|\lambda|-|\mu|},&\mu\prec \lambda,\\
		0,&\text{otherwise}.
	\end{cases}
\end{equation*}

Later we will also use the dual Pieri coefficients
$\psi_{\lambda/\mu}'$ which are defined as
(see \cite[formula~(VI.6.24.iv)]{Macdonald1995} and
\cite[formula~(2.12)]{BorodinCorwin2011Macdonald})
\begin{equation}
	\label{eq:psi_prime_Pieri}
	\psi_{\lambda/\mu}'=
	\psi_{\lambda/\mu}'(q,t)
	\coloneqq
	\prod_{\substack{i<j \\ \lambda_{i}=\mu_{i},\, \lambda_{j}=\mu_{j}+1}}
	\frac{\big(1-q^{\mu_{i}-\mu_{j}} t^{j-i-1}\big)\big(1-q^{\lambda_{i}-\lambda_{j}} t^{j-i+1}\big)}
	{\big(1-q^{\mu_{i}-\mu_{j}} t^{j-i}\big)\big(1-q^{\lambda_{i}-\lambda_{j}} t^{j-i}\big)}.
\end{equation}
for partitions $\lambda$, $\mu$ whose column lengths interlace,
that is,
$\lambda_1'\ge \mu_1'\ge \lambda_2'\ge \mu_2'\ge \cdots $.
We have
\begin{equation}
	\label{eq:P_skew_single_variable_dual}
	P_{\lambda'/\mu'}(x_1 \mid t,q)=
	\begin{cases}
		\psi'_{\lambda/\mu}(q,t)\hspace{1pt} x_1^{|\lambda|-|\mu|}
		,&\mu'\prec \lambda',\\
		0
		,&\text{otherwise}.
	\end{cases}
\end{equation}

\subsection{Cauchy identity}\label{sub:Cauchy_and_Dyson1}

Along with the branching rule, another fundamental
identity for Macdonald polynomials is the Cauchy identity.
Here we present its dual version, see \cite[Chapter~VI.4]{Macdonald1995} for the
usual version.

\begin{Proposition}[{\cite[formula~(VI.5.4) and Chapter~VI.7]{Macdonald1995}}]
	\label{prop:dual_Cauchy_identity}
	Let $N,M\ge 1$ be fixed. We have
	\begin{equation*}%\label{eq:dual_Cauchy}
		\sum_{\lambda\in \mathbb{Y}(N)\colon \lambda_1\le M}
		P_\lambda(x_1,\dots,x_N \mid q,t)
		\hspace{1pt}
		P_{\lambda'}(y_1,\dots,y_M\mid t,q )=
		\prod_{i=1}^{N}\prod_{j=1}^{M}(1+x_iy_j).
	\end{equation*}
	Moreover, for any $\mu\in \mathbb{Y}(N)$ we have
	the following particular case of
	the skew Cauchy identity:
	\begin{equation}		\label{eq:dual_Cauchy_skew}
		P_\mu(x_1,\dots,x_N\mid q,t )
		\prod_{i=1}^{N}(1+x_iy)=
		\sum_{\lambda\in \mathbb{Y}(N)\colon \mu'\prec \lambda'}
		P_{\lambda'/\mu'}(y\mid t,q)
		\hspace{1pt}
		P_\lambda(x_1,\dots,x_N\mid q,t ).
	\end{equation}
\end{Proposition}
Identities~\eqref{eq:branching} and~\eqref{eq:dual_Cauchy_skew} look very similar, but note that
in the former we sum over the smaller partition, while in the latter one we sum over the larger partition.

\subsection{Markov kernels from Macdonald polynomials}\label{eq:Markov_kernels_from_Cauchy}

Let us first recall a general definition from
\cite[Chapter~7]{borodin2016representations}.
Let $\mathfrak{X}$, $\mathfrak{Y}$ be
finite or countable sets. By a Markov kernel
(or a link) from $\mathfrak{X}$ to $\mathfrak{Y}$,
we mean a function $P$ on $\mathfrak{X}\times \mathfrak{Y}$
such that
$P(x,y)\in[0,1]$ for all $x\in \mathfrak{X}$, $y\in \mathfrak{Y}$,
and
\begin{equation*}
	\sum_{y\in \mathfrak{Y}}P(x,y)=1
	\qquad
	\text{for all $x\in \mathfrak{X}$}.
\end{equation*}
We adopt the notation $P\colon \mathfrak{X}\dashrightarrow \mathfrak{Y}$.

Normalizing identities
\eqref{eq:branching} and \eqref{eq:dual_Cauchy_skew}
with nonnegative variables $x_i$ and $y$
(following
\cite{Borodin2010Schur,BorFerr2008DF} in the Schur case and~\cite{BorodinCorwin2011Macdonald}
in the general Macdonald case)
leads to the following two families of links:
\begin{align}
\Lambda^{N}_{N-1}\colon \ & \mathbb{Y}(N)\dashrightarrow \mathbb{Y}(N-1),
		\nonumber\\
		& \Lambda^{N}_{N-1}(\lambda,\mu)\coloneqq
		\psi_{\lambda/\mu}(q,t)\hspace{1pt} x_N^{|\lambda|-|\mu|}\hspace{1pt}
		\frac{
		P_\mu(x_1,\dots,x_{N-1} \mid q,t )
		}
		{P_\lambda(x_1,\dots,x_N \mid q,t)}
		\,\mathbf{1}_{\mu\prec\lambda},\label{eq:traditional_links_Lambda}\\
Q_N\colon \ & \mathbb{Y}(N)\dashrightarrow \mathbb{Y}(N),
		\nonumber\\
		& Q_N(\lambda,\nu)\coloneqq
		\frac{\psi'_{\nu/\lambda}(q,t)
		\hspace{1pt}
		y^{|\nu|-|\lambda|}
		}
		{\prod_{i=1}^{N}(1+x_i y)}
		\frac
		{
		P_\nu(x_1,\dots,x_N\mid q,t )}
		{P_\lambda(x_1,\dots,x_N\mid q,t )}
		\,\mathbf{1}_{\lambda'\prec\nu'}		.\label{eq:traditional_links_Q}
\end{align}
Here and
throughout the paper by $\mathbf{1}_{A}$ we denote
the indicator of the event (or condition) $A$.
For the link \eqref{eq:traditional_links_Q}
we also used \eqref{eq:P_skew_single_variable_dual}.
